% sci_sim_layers.tex — Three simulation layers for the SCI tutorial (S5)
% シミュレーションの3層構造図。SCIチュートリアル S5 用。
%
% Source: docs/architecture/simulation-policy.md §2 (three-layer table),
% §2 "層の関係" (layer 2 is independently positioned, NOT a simplified
% version of layer 3), and §4 (model-match gate).
% 出典: docs/architecture/simulation-policy.md §2（3層構造の表）、
% §2「層の関係」（層2は層3の簡略版ではなく独立した位置づけ）、
% §4（モデル一致ゲート）。
%
% Author's 2nd-round review (2026-09): the previous version drew a FLOW
% (three boxes side by side, left-to-right arrow chain) — the reviewer's
% point was that this reads as a pipeline/flowchart, not as layers. This
% version redraws the three tiers as STACKED HORIZONTAL BANDS (layer1 on
% top, matching the top-to-bottom row order used by the "対応表"/"table"
% frame elsewhere in this chapter), each band self-contained (title / what
% it provides / tool / what it verifies). Short vertical arrows between
% adjacent bands carry the teaching narrative from S5 item 1:
%   モデルがある(層1) -> 実ログで検証できる(層2) -> 飛ばす前にSILSで検証できる(層3)
% The model-match gate is drawn as a SIDE box (not another stacked band,
% since it is not itself a simulation tier — it is the §4 procedure that
% cross-checks layer3's identified plant against layer1's real-machine
% reference), connected by right-angle arrows to layer1 and layer3 only.
%
% 著者の第2ラウンドレビュー（2026-09）: 旧版は3つの箱を横に並べ矢印で
% つなぐ「フロー」だった — 指摘は「層ではなくフローに見える」というもの。
% 本版は3層を縦積みの横長バンドとして描き直す（層1を上に、本章の
% 「対応表」フレームの行順=層1→層2→層3の上から下という並びと揃える）。
% 各バンドは自己完結（タイトル/提供するもの/ツール/何を検証するか）。
% 隣接バンド間の短い垂直矢印が、S5 item 1 の説明の流れを運ぶ:
%   モデルがある(層1) -> 実ログで検証できる(層2) -> 飛ばす前にSILSで検証できる(層3)
% モデル一致ゲートは「層」ではない（§4の判定手順であり、層3で同定した
% プラントを層1の実機基準値と突き合わせる）ため、積み重ねた4つ目のバンドに
% はせず、層1・層3の2つだけに直角矢印で接続するサイド箱として描く。
\documentclass[border=10pt]{standalone}
\usepackage{tikz}
\usetikzlibrary{arrows.meta, positioning, calc}
\usepackage{luatexja}
\usepackage[match]{luatexja-fontspec}
% Cross-platform font: Hiragino (macOS) → Yu Gothic (Windows) → Noto (Linux)
% クロスプラットフォーム: ヒラギノ (macOS) → 游ゴシック (Windows) → Noto (Linux)
\usepackage{ifthen}
\newboolean{jfontset}\setboolean{jfontset}{false}
\IfFontExistsTF{Hiragino Kaku Gothic ProN}{%
  \setmainjfont{Hiragino Kaku Gothic ProN}%
  \setsansjfont{Hiragino Kaku Gothic ProN}%
  \setboolean{jfontset}{true}}{}
\ifthenelse{\boolean{jfontset}}{}{%
  \IfFontExistsTF{Yu Gothic}{%
    \setmainjfont{Yu Gothic}%
    \setsansjfont{Yu Gothic}%
    \setboolean{jfontset}{true}}{}}
\ifthenelse{\boolean{jfontset}}{}{%
  \IfFontExistsTF{Noto Sans CJK JP}{%
    \setmainjfont{Noto Sans CJK JP}%
    \setsansjfont{Noto Sans CJK JP}}{}}

\begin{document}

\definecolor{s1color}{RGB}{0, 100, 180}
\definecolor{s2color}{RGB}{40, 140, 60}
\definecolor{s3color}{RGB}{150, 90, 20}
\definecolor{gatecolor}{RGB}{170, 30, 30}
\definecolor{gray1}{RGB}{120, 120, 120}

% Plain-LaTeX helper macros for the colored text lines inside each band
% 帯の中の色付きテキスト用マクロ
\newcommand{\ttier}[2]{{\bfseries\normalsize\color{#1}#2}}
\newcommand{\tsub}[1]{{\scriptsize\color{black!55}#1}}

\begin{tikzpicture}[
    band/.style={
      rectangle, rounded corners=4pt, draw=#1, line width=1.5pt,
      fill=#1!8, minimum width=108mm, minimum height=13mm,
      align=left, inner xsep=4mm, inner ysep=0.8mm
    },
    gatebox/.style={
      rectangle, rounded corners=4pt, draw=gatecolor, line width=1.3pt,
      fill=gatecolor!7, minimum width=50mm, minimum height=28mm,
      align=left, inner xsep=3mm, inner ysep=2mm
    },
    relarr/.style={-{Stealth[length=3mm]}, line width=1.5pt, gray1},
    garr/.style={-{Stealth[length=3mm]}, line width=1.4pt, gatecolor},
  ]

  % Three STACKED bands, top (Layer 1, design-time) to bottom (Layer 3,
  % SILS) — matches the row order of the "対応表" table elsewhere in this
  % chapter (層1/層2/層3 top-to-bottom). This is the core fix for the
  % reviewer's "flow, not layers" complaint: full-width horizontal bars
  % stacked vertically, not boxes chained left-to-right.
  %
  % Line budget: each band is title + 2 SHORT sublines (not long sentences)
  % so every line fits in ONE row at this width without wrapping — a
  % wrapped line silently doubles a band's height and was the actual cause
  % of an earlier revision's 91pt (32mm) frame overflow (verified by
  % rebuilding and reading the log, not guessed).
  % 3つの帯を縦に積む。上（層1・設計時）から下（層3・SILS）へ — 本章の
  % 「対応表」の行順（層1/層2/層3、上から下）と揃える。「フローであって
  % 層ではない」という指摘への核心的な修正: 横長のバーを縦に積むのであって、
  % 箱を左から右へ連鎖させるのではない。
  %
  % 行数の設計: 各帯はタイトル＋短い2行（長い一文にしない）とし、この幅で
  % 折り返さず1行に収める — 行の折り返しは帯の高さを暗黙に倍増させ、
  % 以前の版でフレームが91pt（32mm）溢れた実際の原因だった（推測ではなく
  % 再ビルドしてログを読んで確認済み）。
  \node[band=s1color] (S1) at (0, 0) {%
    \begin{minipage}{99mm}\raggedright
      \ttier{s1color}{層1: 設計用線形モデル}\\[1.5pt]
      \tsub{$G(s){=}b\,e^{-Ls}/(s(Ts{+}1))$（伝達関数として扱う）}\\[1pt]
      \tsub{\texttt{sf sysid rate-fit} で同定 $\to$ 設計・ループ整形の正}
    \end{minipage}%
  };
  \node[band=s2color, below=5mm of S1] (S2) {%
    \begin{minipage}{99mm}\raggedright
      \ttier{s2color}{層2: 実ログ駆動オフライン再生}\\[1.5pt]
      \tsub{層1と同じ$G(s)$＋飽和等，制御則はPython移植}\\[1pt]
      \tsub{実ログ外乱で駆動（\texttt{analysis/scripts/}）$\to$ A/B判定の正}
    \end{minipage}%
  };
  \node[band=s3color, below=5mm of S2] (S3) {%
    \begin{minipage}{99mm}\raggedright
      \ttier{s3color}{層3: SILS（実ファーム + MuJoCo）}\\[1.5pt]
      \tsub{非線形6自由度のプラント＋実ファームC++そのもの}\\[1pt]
      \tsub{Code/Param Identity $\to$ 飛ばす前のシステム全体検証}
    \end{minipage}%
  };

  % Short vertical arrows between adjacent bands, straight (no bend needed
  % since the bands already share the same x-center), centered on the gap
  % between bands. Each carries one clause of the S5 item-1 narrative.
  % 隣接する帯の間の短い垂直矢印。帯は既に同じx中心を共有しているため
  % 折れ曲がり不要の直線。帯間のギャップの中央に配置し、各矢印がS5 item 1
  % の説明文の一節を運ぶ。
  \draw[relarr] (S1.south) -- (S2.north);
  \draw[relarr] (S2.south) -- (S3.north);

  % Labels sit to the RIGHT of each vertical arrow (not on top of the
  % line — T2), inside the inter-band gap so they cannot collide with any
  % band's text (the gap is empty space between two boxes) or with the
  % gate/bus on the right margin (labels stop well before that region).
  % ラベルは各垂直矢印の右側に置く（線の上に乗せない — T2）。帯間ギャップ
  % （2つの箱の間の空白）の中に収まるため、帯のテキストとも右マージンの
  % ゲート/バスとも衝突しない（ラベルはその領域よりずっと手前で止まる）。
  \node[font=\scriptsize, text=gray1, anchor=west]
    at ($(S1.south)!0.5!(S2.north)+(4mm,0)$) {実ログで照合};
  \node[font=\scriptsize, text=gray1, anchor=west]
    at ($(S2.south)!0.5!(S3.north)+(4mm,0)$) {飛行前検証};

  % Model-match gate (§4): a SIDE box, not a fourth stacked layer — it is
  % the verification procedure that cross-checks layer3's identified
  % plant against layer1's real-machine reference, so it connects to
  % layer1 AND layer3 only (never layer2, which the policy explicitly
  % keeps independent of layer3 — §2 "層の関係"). Vertically centered on
  % the S1-S3 span (== S2's height, since bands are symmetric).
  % モデル一致ゲート（§4）: 「層」ではなくサイド箱 — 層3で同定したプラント
  % を層1の実機基準値と突き合わせる判定手順そのものなので、層1・層3にのみ
  % 接続する（層2には繋がない。方針書§2「層の関係」で層2は層3と独立と
  % 明記されているため）。S1〜S3の全体スパンの中央（帯が対称なのでS2の高さ
  % と同じ）に配置。
  \node[gatebox, right=20mm of S2] (GATE) {%
    \begin{minipage}{43mm}\raggedright
      \ttier{gatecolor}{モデル一致の合否判定}\\[3pt]
      \tsub{層3ログを層1と同一手順で同定}\\[2pt]
      \tsub{$(b,L,T)$ を層1基準値と比較}\\[2pt]
      \tsub{（\texttt{sf sils sysid-gate}）}
    \end{minipage}%
  };

  % Right-angle routing from Layer 1 / Layer 3 into the gate. The corner
  % x-coordinate is taken from GATE (via "-|", which reads x from its
  % RIGHT operand) so the vertical run lands exactly on GATE's north/south
  % edge midpoint — an edge, not a corner (T4f) — with a >=8mm straight
  % trunk before each arrowhead.
  % 層1／層3からゲートへの直角配線。コーナーのx座標はGATE側から取る
  % （"-|" は右側の被演算子からxを読む）ため、垂直区間はGATEのnorth/south
  % 辺の中点に正確に着地する（角ではなく辺 — T4f）。矢頭前に8mm以上の
  % 直線区間を確保する。
  \draw[garr] (S1.east) -- (S1.east -| GATE) -- (GATE.north);
  \draw[garr] (S3.east) -- (S3.east -| GATE) -- (GATE.south);

  % Labels for the gate's two inputs, placed above/below their horizontal
  % run (clear of the S1/S3 bands to the left and the GATE box to the
  % right — this horizontal segment is empty space by construction).
  % ゲートへの2入力のラベル。それぞれの水平区間の上/下に配置する
  % （左のS1/S3帯とも右のGATE箱とも構造上ぶつからない空白領域）。
  \node[font=\scriptsize, text=gatecolor, anchor=south]
    at ($(S1.east)!0.5!(S1.east -| GATE)+(0,1mm)$) {実機同定値（基準）};
  \node[font=\scriptsize, text=gatecolor, anchor=north]
    at ($(S3.east)!0.5!(S3.east -| GATE)+(0,-1mm)$) {同一パイプラインで同定};

\end{tikzpicture}
\end{document}
